% short/01_introduction.tex: the question, why it is hard, the approach, the status, related work (about 1 page
% with Figures 1 and 2). Full version: sections/01_introduction.tex. The status paragraph is the ONLY place that
% states the audit status of the final-model lane.
\section{Introduction}\label{sec:intro}

\paragraph{The question.} A person holds a smartphone at arm's length and brings it to the nose with the left or
the right hand (Figure~\ref{fig:gesture}). The task, set by the principal investigator (PI), is to predict which
hand from the motion sensors that the browser's DeviceMotion interface delivers at about
60~Hz~\cite{w3c2025deviceorientation}, as a calibrated probability \pright{} that the right hand moved the phone
(of the touches given 80\,\%, about 80\,\% should be right-hand touches), for persons the model has never seen.

% Figures 1 and 2 share one float (the gesture frames stacked beside the diagram), placed here so that it lands at
% the top of the next page and the introduction flows on below it.
\begin{figure}[tbp]
\centering
\begin{minipage}[c]{0.34\linewidth}\centering
  \begin{subfigure}{\linewidth}\centering
    \repoimage{gesture-a}{../docs/img/gesture/right_upright_natural_start.png}\caption{Arm extended, still}
  \end{subfigure}\\[3pt]
  \begin{subfigure}{\linewidth}\centering
    \repoimage{gesture-b}{../docs/img/gesture/right_upright_natural_halfway.png}\caption{The approach}
  \end{subfigure}\\[3pt]
  \begin{subfigure}{\linewidth}\centering
    \repoimage{gesture-c}{../docs/img/gesture/right_upright_natural_touch.png}\caption{Nose touch and hold}
  \end{subfigure}
\end{minipage}\hfill
\begin{minipage}[c]{0.58\linewidth}\centering
  % short/diagram_pipeline.tex: the overview diagram, shared by short.tex and brief.tex (\input it inside a figure).
% TikZ, drawn at a fixed size (about 7.4 cm wide, explicit 7.5/8.5-pt fonts), so it reads the same at column width in the
% brief and beside its caption in the short version. Numbers only through \R keys. Dashed boxes: the cue-free steps,
% simulated by lane m4 with its reference classifier yaw8; the final model has not yet run behind the detector.
\begin{tikzpicture}[
    node distance=2.4mm,
    stage/.style={draw=black!55, rounded corners=1.5pt, fill=black!3, text width=6.8cm, inner sep=2.4pt,
      align=left, font=\fontsize{7.5}{9}\selectfont},
    cuefree/.style={stage, dashed, fill=white},
    flow/.style={-{Latex[length=1.5mm,width=1.3mm]}, black!65, semithick},
    head/.style={font=\fontsize{8.5}{10}\selectfont\bfseries}]
  \node[stage] (sensors) {{\fontsize{8.5}{10}\selectfont\bfseries 1\enspace Phone sensors}\\
    accelerometer and gyroscope, read by the browser at about 60~Hz};
  \node[cuefree, below=of sensors] (detect) {{\fontsize{8.5}{10}\selectfont\bfseries 2\enspace Find the approach
    (no start cue)}\\ cut the stream into candidate movements; take the most likely approach if
    $P(\mathrm{approach}) \geq \R{m4/det_threshold}$, else wait};
  \node[stage, below=of detect] (std) {{\fontsize{8.5}{10}\selectfont\bfseries 3\enspace Standardise}\\
    gravity frame and heading-following body frame; \R{p1/n_norm} samples, 12 channels};
  \node[stage, below=of std] (model) {{\fontsize{8.5}{10}\selectfont\bfseries 4\enspace Model}\\
    1D CNN, trained only on other persons};
  \node[stage, below=of model] (calib) {{\fontsize{8.5}{10}\selectfont\bfseries 5\enspace Calibrated
    \pright{}}\\ Platt calibrator fitted inside each training fold};
  \node[cuefree, below=of calib] (stop) {{\fontsize{8.5}{10}\selectfont\bfseries 6\enspace Several attempts,
    then stop}\\ combine the logits, $\sigma\bigl(\sum_i z_i/\tau_k\bigr)$; answer once the result leaves
    $(1 - p_{\mathrm{stop}},\, p_{\mathrm{stop}})$, $p_{\mathrm{stop}} = \R{m4/p_stop}$, or after
    \R{m4/k_max} attempts\\[13pt]\mbox{}};
  % The answer bar inside box 6 (schematic, not to scale): left of the band LEFT, right of it RIGHT.
  \begin{scope}[shift={(stop.south west)}, xshift=2.5mm, yshift=1.4mm, font=\fontsize{7.5}{9}\selectfont]
    \fill[handleft!30]  (0,0) rectangle (1.25,0.30);
    \fill[black!10]     (1.25,0) rectangle (5.05,0.30);
    \fill[handright!30] (5.05,0) rectangle (6.3,0.30);
    \draw[black!50] (0,0) rectangle (6.3,0.30);
    \node[text=handleft]  at (0.625,0.15) {\sffamily\bfseries LEFT};
    \node[text=black!65]  at (3.15,0.15) {next attempt};
    \node[text=handright] at (5.675,0.15) {\sffamily\bfseries RIGHT};
  \end{scope}
  \foreach \a/\b in {sensors/detect, detect/std, std/model, model/calib, calib/stop}
    \draw[flow] (\a) -- (\b);
  \draw[flow] (stop.east) -- ++(3.2mm,0) |- (detect.east); % back to step 2 while P stays inside the band
  \node[below=1.2mm of stop.south west, anchor=north west, inner sep=0pt, font=\fontsize{7}{8.5}\selectfont,
    text=black!70] {dashed: simulated with lane m4's classifier yaw8};
\end{tikzpicture}%

\end{minipage}\par\medskip
\captionsetup{format=plain,justification=raggedright,singlelinecheck=false}
\begin{minipage}[t]{0.34\linewidth}
  \caption{The gesture as the app shows it (right hand, upright grip, natural speed). The model reads the
    approach.}
  \label{fig:gesture}
\end{minipage}\hfill
\begin{minipage}[t]{0.58\linewidth}
  \caption{From sensors to an answer. Steps 3--5 are evaluated on the stored approach windows by leaving one person
    out; the dashed cue-free steps were simulated with lane m4's reference classifier, not yet with the final
    model.}
  \label{fig:pipeline}
\end{minipage}
\end{figure}

\paragraph{Why it is hard.} Persons, phones and grips differ, and the device axes turn with the grip. The two hands
make mirror-image movements, so most simple descriptors (speed, duration, path length) are the same for both, and
in production there is no start cue: the page must find the approach in a continuous stream.

\paragraph{The approach.} The left--right mirror symmetry of the body shows which quantities can carry the hand
(Section~\ref{sec:physics}). Each trial is standardised into a gravity frame and a heading-following body frame, four
model families are compared by leave-one-person-out (LOSO) cross-validation, a pre-specified rule picks the final
model, and a simulation then removes the start cue (Figure~\ref{fig:pipeline}).

\paragraph{Status.} This preliminary version shortens the full report (report/main.pdf in
\url{https://github.com/AxelFridman/nosetouch-study}), which states every decision with its evidence. The analysis ran in lanes: the three model-family lanes (m1--m3), preprocessing (p1) and the
production simulation (m4) each passed two independent audits; the final-model lane is a first pass, not yet
independently audited, and the final model has not yet run in the cue-free simulation.

\paragraph{Related work.} Most prior studies infer the hand from touch traces or grip during ordinary phone use, and
several split the data by trial or at random, which overstates the accuracy for new
persons~\cite{goel2012gripsense,guo2016operating,fernandez2021implicit,saeb2017usecase,little2017crossvalidation}.
Opposite rotation senses of the two hands were seen with wrist sensors in a clinical finger-to-nose test, without a
classifier~\cite{mohammadighazi2022ataxia}, and modern classifiers are often over-confident~\cite{guo2017calibration},
so this study evaluates person-wise and checks calibration.
